# Big Data Banking Analytics Pipeline

## 📊 Overview
This project builds an end-to-end big data architecture for analyzing customer spending behavior in a banking environment. The system integrates relational databases, data warehousing, NoSQL storage, and a front-end analytics interface.

The goal is to enable scalable, fast, and flexible analysis of customer transactions for insights into spending patterns, customer segmentation, and decision-making.

---

## 🧠 Architecture Summary
The system is designed using a modern data architecture with multiple layers:

1. **Data Source & OLTP Layer (MySQL RDS)**
2. **Data Warehouse Layer (Hive on AWS EMR)**
3. **Speed Layer (DynamoDB NoSQL)**
4. **Analytics Layer (Python)**
5. **Front-End Interface**

---

## 🗄️ Data Engineering (Relational System)
A synthetic banking database was constructed using real-world datasets:

- ~54,000 customers
- ~120,000 accounts
- ~30,000 loans
- ~280,000 transactions

### Tables:
- `customer`
- `account`
- `loan`
- `transaction_table`

The schema enforces:
- Primary and foreign key relationships
- Cleaned and standardized attributes
- Referential integrity validation

---

## ⚡ NoSQL Speed Layer (DynamoDB)
To support fast queries, a pre-aggregated table was built:

- Table: `cust_aggregates`
- Key: `(customer_id, month)`
- Includes:
  - Total monthly spending
  - Category-level summaries
  - Overspending indicators

### Why DynamoDB?
- Low-latency reads
- Scalable for real-time applications
- Avoids repeated heavy queries on transactional data

---

## 🏗️ Data Warehouse (Hive on EMR)
A star schema was implemented for analytics:

### Fact Table:
- `fact_transaction`

### Dimension Tables:
- `dim_customer`
- `dim_account`
- `dim_category`
- `dim_date`

### Key Features:
- Partitioning by year and month
- Efficient time-based queries
- Denormalized analytics table:
  - `customer_spend_analytics`

---

## 📈 Analytics Pipeline (Python + Hive)
A Python-based analytics tool was developed to:

- Query Hive using `pyhive`
- Filter data by user-defined inputs (e.g., age range)
- Compute:
  - Total and average spending
  - Category-level spending
  - Demographic breakdowns

### Outputs:
- Text-based reports
- Multi-panel visualization (PNG)
  - Spending by category
  - Spending by job, education, marital status
  - Category distribution

---

## 🖥️ Front-End Application
An interactive interface allows users to:

- Input parameters (e.g., age range)
- Trigger backend analytics on EMR
- View generated visualizations

The front end communicates with the EMR cluster, executes Python + Hive queries, and returns results dynamically.

---

## 🛠️ Tech Stack
- Python (Pandas, PyHive)
- MySQL (RDS)
- Apache Hive (EMR)
- Hadoop (HDFS)
- AWS DynamoDB
- AWS S3
- Flask (lightweight front end)

---

## 🚀 Key Takeaways
- Built a full data pipeline from raw data → analytics → visualization
- Demonstrated understanding of:
  - OLTP vs OLAP systems
  - Star schema design
  - Distributed data processing (Hive/EMR)
  - NoSQL optimization for speed
- Designed a system that balances performance, scalability, and usability

